\uuid{hqCN}
\titre{Théorème de Rolle}
\auteur{bodin ; hanani ; mzari}
\organisation{exo7}
\licence{CC BY-NC-SA 4.0}
\datecreate{2020-01-20}
\niveau{L1}
\module{Analyse}
\chapitre{Dérivabilité des fonctions réelles}
\sousChapitre{Autre}
\difficulte{3}
\theme{Dérivées, théorème de Rolle}
\exo7qcmid{lille-1.derivees.024}

\contenu{
\begin{question}
\qclassification{124.00}
\qtags{motcle=théorème de Rolle}
Soit $f(x)=\tan (x)$. Quelles sont les bonnes réponses ?
\begin{answers}
    \bad{$f(0)=0=f(\pi)$ et donc il existe $c\in ]0,\pi[$ tel que $f'(c)=0$.}
    \good{$f(0)=0=f(\pi)$ mais il n'existe pas de $c\in ]0,\pi[$ tel que $f'(c)=0$.}
    \bad{Le théorème de Rolle ne s'applique pas à $f$ sur $[0,\pi]$ car $f(0)\neq f(\pi)$.}
    \good{Le théorème de Rolle ne s'applique pas à $f$ sur $[0,\pi]$.}
\end{answers}
\begin{explanations}
On a bien $f(0)=0=f(\pi)$. Mais, pour tout $\displaystyle x\neq \frac{\pi}{2}+k\pi$, $k\in \Zz$, on a $f'(x)=1+\tan ^2x>0$. On ne peut appliquer le théorème de Rolle à $f$ sur $[0,\pi]$ car $f$ n'est pas définie au point $\displaystyle \frac{\pi}{2}$.
\end{explanations}
\end{question}
}
